\documentclass[11pt,a4paper]{article}

% --- Pakker ---
\usepackage[utf8]{inputenc}
\usepackage[T1]{fontenc}
\usepackage[danish,provide=*]{babel}
\usepackage{amsmath, amssymb}
\usepackage[a4paper, margin=2.2cm]{geometry}
\usepackage{tikz}
\usepackage{pgfplots}
\pgfplotsset{compat=1.18}
\usetikzlibrary{arrows.meta,decorations.markings,calc,patterns}
\usepackage{booktabs}
\usepackage{enumitem}
\usepackage{xcolor}
\usepackage{wasysym}
\usepackage{hyperref}

\hypersetup{
  colorlinks=true,
  linkcolor=blue!50!black,
  urlcolor=blue!50!black
}

% Farver
\definecolor{bgpurple}{HTML}{7C5CFC}
\definecolor{bgcyan}{HTML}{38BDF8}
\definecolor{bgink}{HTML}{0B0D17}
\definecolor{bgink50}{HTML}{7A7E8E}
\definecolor{bggold}{HTML}{F0C419}

% --- Opgave-makro ---
\newcommand{\opg}[2]{%
  \par\vspace{0.6em}\noindent
  \textbf{\large #1}\par
  #2\par
}

% --- Facit-makro (kompakt, inline efter opgave) ---
\newcommand{\facit}[1]{%
  \par\vspace{0.3em}\noindent
  \begin{tikzpicture}
    \node[draw=bggold!80, fill=bggold!8, rounded corners=3pt, inner sep=6pt,
          text width=\linewidth-2\fboxsep-2\pgflinewidth-9pt, align=left]
          {\small\textbf{Facit:}\ #1};
  \end{tikzpicture}\par
}

% --- Regel-boks (forklaring) ---
\newcommand{\regel}[2]{%
  \par\vspace{0.4em}\noindent
  \begin{tikzpicture}
    \node[draw=bgcyan!70, fill=bgcyan!10, inner sep=8pt,
          text width=\linewidth-2\fboxsep-2\pgflinewidth-9pt, align=left]
          {\textbf{Regel:} #1\\[2pt]\textit{#2}};
  \end{tikzpicture}\par
}

% --- Eksempel-boks ---
\newcommand{\eksempel}[2]{%
  \par\vspace{0.3em}\noindent
  \fcolorbox{bgink!30}{yellow!6}{%
    \parbox{0.96\linewidth}{%
      \textbf{Eksempel:}\ \textit{#1}\\[3pt]
      #2
    }%
  }\par
}

% --- Crash-boks ---
\newcommand{\crash}[1]{%
  \par\vspace{0.3em}\noindent
  \begin{tikzpicture}
    \node[draw=bgpurple!80, fill=bgink!3, rounded corners=4pt, inner sep=8pt,
          text width=\linewidth-2\fboxsep-2\pgflinewidth-9pt, align=left]
          {\textbf{\color{bgpurple}Crash-forklaring:}\\[2pt]#1};
  \end{tikzpicture}\par
}

% --- Lærer-tip (separat boks) ---
\newcommand{\laerertip}[1]{%
  \par\vspace{0.2em}\noindent
  \begin{tikzpicture}
    \node[draw=bgink!40, fill=bgink!4, inner sep=5pt,
          text width=\linewidth-2\fboxsep-2\pgflinewidth-9pt, align=left]
          {\footnotesize\textbf{\color{bgink}Tip til læreren:}\ #1};
  \end{tikzpicture}\par
}

\title{\textbf{GXU Lærervejledning 4}\\[0.3em]
      \large Brøkregning, Statistik \& Trigonometri\\
      \large\textit{(med facit mellem hver opgave)}}
\author{Fra Omar og Karoline \smiley}
\date{Oktober 2026}

\begin{document}

\maketitle
\noindent\textbf{Formål:} Lærer-udgave af GXU Træningsopgaver 4. Hver opgave har en gul \textbf{Facit-boks} lige under sig, så du kan rette mens du kører rundt. Kildehenvisninger er samlet på side~1. Plot-eksempler og teori er de samme som elev-hæftet — kun facit er flyttet fra appendiks til umiddelbart efter hver opgave.

\vspace{0.5em}\noindent\rule{\linewidth}{0.4pt}

\begin{center}
\footnotesize\color{bgink!90}
\textbf{Kilder og inspiration:} Webmatematik (\textit{webmatematik.dk}, Matematikcenter) · Khan Academy · Matematik i grundforløbet (emu.dk) · Geogebra.org.\\[2pt]
Eksempler med \textit{(kilde: ...)} bygger på disse sider — de er genfortalt og tilpasset, men idé og metode kommer derfra. Crash-forklaringerne er vores egne.
\end{center}
\vspace{0.3em}\noindent\rule{\linewidth}{0.4pt}

% ===========================================================
\section{Brøkregning — forlæng, forkort, addér og gange}
% ===========================================================

\crash{Brøker handler om \textbf{andele}. Tælleren (øverst) fortæller hvor mange dele vi har, nævneren (nederst) fortæller hvor mange lige store dele det hele er delt i. Når vi regner med brøker, er den vigtigste regel: \emph{vi må frit forlænge og forkorte} — det ændrer ikke værdien af brøken. Forestil dig en chokoladeplade delt i 4 ruder: 2/4 er præcis det samme som 1/2.}

\noindent To grundregler for brøker med samme nævner:
\begin{itemize}[leftmargin=2em,itemsep=2pt]
\item \textbf{Addition:} $\dfrac{a}{n}+\dfrac{b}{n}=\dfrac{a+b}{n}$ — læg tællere sammen, behold nævner.
\item \textbf{Subtraktion:} $\dfrac{a}{n}-\dfrac{b}{n}=\dfrac{a-b}{n}$ — træk tællere fra hinanden, behold nævner.
\end{itemize}

\noindent For brøker med \emph{forskellige} nævner skal vi først finde en \textbf{fællesnævner} — typisk ved at gange nævnerne med hinanden, eller bedre: find den \textbf{mindste fællesnævner} (mindste fælles multiplum).

\regel{For at addere/subtrahere brøker: find en fællesnævner, forlæng hver brøk så de matcher, og læg så tællerne sammen.}{%
\textbf{Husk:} forkort altid til sidst, så svaret står i simplest form.}

\eksempel{Forlæng og forkort $\dfrac{2}{5}$}{%
  Forlæng med 3: $\dfrac{2}{5}=\dfrac{2\cdot 3}{5\cdot 3}=\dfrac{6}{15}$.\\[3pt]
  Forkort $\dfrac{12}{20}$ med 4: $\dfrac{12}{20}=\dfrac{12:4}{20:4}=\dfrac{3}{5}$.\\[3pt]
  \textit{Fortolkning:} vi ganger (eller dividerer) både tæller og nævner med det samme tal — brøkens værdi ændrer sig ikke, kun måden vi skriver den på.
}

\eksempel{Addér $\dfrac{1}{3}+\dfrac{1}{4}$}{%
  Fællesnævner for 3 og 4 er 12 (mindste fælles multiplum).\\[3pt]
  Forlæng: $\dfrac{1}{3}=\dfrac{4}{12}$ og $\dfrac{1}{4}=\dfrac{3}{12}$.\\[3pt]
  Læg sammen: $\dfrac{4}{12}+\dfrac{3}{12}=\dfrac{7}{12}$.\\[3pt]
  \textit{Tjek:} 7/12 kan ikke forkortes yderligere (7 er et primtal, der ikke går op i 12).
}

\eksempel{Gang to brøker $\dfrac{2}{3}\cdot\dfrac{5}{7}$}{%
  Regel: tæller gange tæller, nævner gange nævner.\\[3pt]
  $\dfrac{2}{3}\cdot\dfrac{5}{7}=\dfrac{2\cdot 5}{3\cdot 7}=\dfrac{10}{21}$.\\[3pt]
  \textit{Bemærk:} 10 og 21 har ingen fælles faktorer, så brøken er i simplest form.
}

\eksempel{Divider $\dfrac{3}{4}:\dfrac{2}{5}$}{%
  Regel: dividér med at \textbf{gange med den omvendte brøk}.\\[3pt]
  $\dfrac{3}{4}:\dfrac{2}{5}=\dfrac{3}{4}\cdot\dfrac{5}{2}=\dfrac{15}{8}=1\tfrac{7}{8}$.\\[3pt]
  \textit{Tjek:} et stort tal divideret med et lille tal giver noget over 1 — det giver mening, for $\tfrac{3}{4}>\tfrac{2}{5}$.
}

\laerertip{Opgaverne 1.1--1.4 dækker hele grundridsen: forkort $\to$ addition med fællesnævner $\to$ multiplikation/division $\to$ anvendelse. 1.4 er ``omvendt proportionalitet'' i forklædning — en god elev skal kunne se at forholdet $6/4$ skalerer brøken. Hvis en elev hænger i 1.2, så lad dem først prøve med ``samme nævner''-eksempler (som $\tfrac{3}{7}+\tfrac{2}{7}$).}

\opg{Opgave 1.1}{Forkort følgende brøker mest muligt: (a) $\dfrac{8}{12}$, (b) $\dfrac{15}{25}$, (c) $\dfrac{18}{30}$.}
\facit{(a) $\tfrac{8}{12}=\tfrac{2}{3}$ (divider med 4). \quad (b) $\tfrac{15}{25}=\tfrac{3}{5}$ (divider med 5). \quad (c) $\tfrac{18}{30}=\tfrac{3}{5}$ (divider med 6).}

\opg{Opgave 1.2}{Regn ud og angiv svaret i simplest form:
\[
\frac{2}{5}+\frac{3}{4}, \qquad \frac{5}{6}-\frac{1}{3}, \qquad \frac{3}{7}+\frac{2}{7}.
\]}
\facit{$\tfrac{2}{5}+\tfrac{3}{4}=\tfrac{8}{20}+\tfrac{15}{20}=\tfrac{23}{20}=1\tfrac{3}{20}$.\\[2pt]
$\tfrac{5}{6}-\tfrac{1}{3}=\tfrac{5}{6}-\tfrac{2}{6}=\tfrac{3}{6}=\tfrac{1}{2}$.\\[2pt]
$\tfrac{3}{7}+\tfrac{2}{7}=\tfrac{5}{7}$ (allerede samme nævner).}

\opg{Opgave 1.3}{Regn ud:
\[
\frac{2}{3}\cdot\frac{5}{8}, \qquad \frac{7}{4}\cdot\frac{3}{14}, \qquad \frac{3}{4}:\frac{5}{2}.
\]}
\facit{$\tfrac{2}{3}\cdot\tfrac{5}{8}=\tfrac{10}{24}=\tfrac{5}{12}$ (divider med 2).\\[2pt]
$\tfrac{7}{4}\cdot\tfrac{3}{14}=\tfrac{21}{56}=\tfrac{3}{8}$ (divider med 7).\\[2pt]
$\tfrac{3}{4}:\tfrac{5}{2}=\tfrac{3}{4}\cdot\tfrac{2}{5}=\tfrac{6}{20}=\tfrac{3}{10}$.}

\opg{Opgave 1.4 (bonus)}{En opskrift på 4 personer bruger $\tfrac{2}{3}$ dl olie. Hvor meget olie skal du bruge til 6 personer, hvis mængden vokser proportionalt?}
\facit{Forholdet mellem personer: $\tfrac{6}{4}=\tfrac{3}{2}=1{,}5$.\\[2pt]
Olie: $1{,}5\cdot\tfrac{2}{3}\text{ dl}=\tfrac{3}{2}\cdot\tfrac{2}{3}\text{ dl}=\tfrac{6}{6}\text{ dl}=1\text{ dl}$.}

% ===========================================================
\section{Statistik og sandsynlighed — middelværdi, median og P}
% ===========================================================

\crash{Statistik handler om at \textbf{opsummere mange tal} til ét eller to nøgletal, man kan sammenligne. Sandsynlighed handler om at \textbf{forudsige chancen} for en hændelse. Begge dele bygger på den samme idé: tæl gunstige udfald og sammenlign med det samlede antal muligheder.}

\noindent Centrale begreber:
\begin{itemize}[leftmargin=2em,itemsep=2pt]
\item \textbf{Stikprøve} (dataset): de tal vi har målt.
\item \textbf{Middelværdi (gennemsnit):} $\bar x = \dfrac{x_1+x_2+\ldots+x_n}{n}$.
\item \textbf{Median:} den midterste værdi når data er sorteret (eller gennemsnittet af de to midterste, hvis $n$ er lige).
\item \textbf{Median vs.\ middelværdi:} medianen er robust over for ekstreme værdier (outliers), middelværdien er følsom.
\item \textbf{Sandsynlighed:} $P(\text{hændelse})=\dfrac{\text{antal gunstige udfald}}{\text{antal mulige udfald}}$, når alle udfald er lige sandsynlige.
\end{itemize}

\regel{Middelværdi: læg alle tal sammen og divider med antal. Median: sortér og tag midten.}{%
\textbf{Huskeregel:} hvis der er outliers (fx én elev med karakteren $-3$ og resten med 2--12), så er \textbf{medianen} et bedre bud på ``det typiske'' end middelværdien.
}

\eksempel{Middelværdi og median for $2, 5, 5, 7, 11$}{%
  Sum: $2+5+5+7+11=30$. Antal tal: $5$.\\[3pt]
  Middelværdi: $\bar x=\dfrac{30}{5}=6$.\\[3pt]
  Median: data er allerede sorteret, den midterste (3.\,største) er $\mathbf{5}$.\\[3pt]
  \textit{Bemærk:} middelværdien (6) er større end medianen (5), fordi 11 trækker op. Medianen er et bedre bud på ``det typiske tal'' her.
}

\eksempel{Sandsynlighed for en bestemt terning}{%
  En almindelig terning har 6 sider, alle lige sandsynlige.\\[3pt]
  $P(\text{slå et 4-tal})=\dfrac{1}{6}\approx 16{,}7\%$.\\[3pt]
  $P(\text{slå et lige tal})=\dfrac{3}{6}=\dfrac{1}{2}=50\%$ (udfald: 2, 4, 6).\\[3pt]
  \textit{Fortolkning:} ``lige tal'' rummer 3 ud af 6 muligheder, så sandsynligheden er 50\%.
}

\eksempel{To terningekast — hvad er $P(\text{to 6'ere})$?}{%
  Kast 1: $P(6)=\tfrac{1}{6}$. Kast 2: $P(6)=\tfrac{1}{6}$ (uafhængige kast).\\[3pt]
  $P(\text{begge}=6)=\tfrac{1}{6}\cdot\tfrac{1}{6}=\tfrac{1}{36}\approx 2{,}8\%$.\\[3pt]
  \textit{Regel:} uafhængige hændelser ganges. ``Uafhængig'' betyder at det ene kast ikke påvirker det andet.
}

\laerertip{2.1--2.4: start med ``tæl-og-divider''-spørgsmål, slut med kombinatorik (terning). 2.4 (bonus) tester om eleven kan generalisere: gul kugle findes ikke, så $P=0$. Det er en god ``fælde''-opgave. 2.3 kræver at eleven enten laver en 6x6-tabel eller husker de symmetriske tællinger (1, 2, 3, 4, 5, 6, 5, 4, 3, 2, 1).}

\opg{Opgave 2.1}{Beregn middelværdi og median for datasættet: $3, 7, 7, 2, 9, 4, 7, 1, 8$.}
\facit{Sorteret: $1, 2, 3, 4, 7, 7, 7, 8, 9$. Sum $=48$, $n=9$.\\[2pt]
Middelværdi: $\bar x=\tfrac{48}{9}\approx 5{,}33$.\\[2pt]
Median: 5.\,værdi (midten) $=7$.\\[2pt]
\textit{Bemærk:} middelværdien er lavere end medianen, fordi 1 og 2 trækker ned.}

\opg{Opgave 2.2}{I en klasse på 24 elever er der 14 piger og 10 drenge. Hvad er sandsynligheden for tilfældigt at udtrække (a) en pige, (b) en dreng, (c) en elev der enten er pige eller dreng?}
\facit{(a) $P(\text{pige})=\tfrac{14}{24}=\tfrac{7}{12}\approx 58{,}3\%$.\\[2pt]
(b) $P(\text{dreng})=\tfrac{10}{24}=\tfrac{5}{12}\approx 41{,}7\%$.\\[2pt]
(c) $P(\text{pige eller dreng})=\tfrac{24}{24}=1$ (sikker hændelse).}

\opg{Opgave 2.3}{Du kaster to terninger. Hvad er sandsynligheden for at summen af de to øjne er (a) 7, (b) højst 4, (c) præcis 12? Tip: lav en tabel over alle 36 udfald.}
\facit{Mulige summer: $2$ (1 komb.), $3$ (2), $4$ (3), $5$ (4), $6$ (5), $7$ (6), $8$ (5), $9$ (4), $10$ (3), $11$ (2), $12$ (1). I alt 36.\\[2pt]
(a) $P(7)=\tfrac{6}{36}=\tfrac{1}{6}\approx 16{,}7\%$.\\[2pt]
(b) $P(\leq 4)=\tfrac{1+2+3}{36}=\tfrac{6}{36}=\tfrac{1}{6}\approx 16{,}7\%$.\\[2pt]
(c) $P(12)=\tfrac{1}{36}\approx 2{,}8\%$.}

\opg{Opgave 2.4 (bonus)}{En urne indeholder 3 røde, 5 blå og 2 grønne kugler. Du trækker én kugle tilfældigt. Hvad er sandsynligheden for (a) rød, (b) ikke-grøn, (c) gul? Hvad er den samlede sandsynlighed for alle tre farver?}
\facit{Total: $3+5+2=10$ kugler.\\[2pt]
(a) $P(\text{rød})=\tfrac{3}{10}=30\%$.\\[2pt]
(b) $P(\text{ikke-grøn})=\tfrac{8}{10}=80\%$.\\[2pt]
(c) $P(\text{gul})=0$ (ingen gule kugler i urnen).\\[2pt]
Total $P$ for de tre farver: $P(\text{rød})+P(\text{blå})+P(\text{grøn})=30\%+50\%+20\%=100\%$ — alle udfald dækket.}

% ===========================================================
\section{Trigonometri — retvinklede trekanter}
% ===========================================================

\crash{Trigonometri handler om \textbf{forholdet mellem siderne} i en retvinklet trekant. Vi definerer tre funktioner — $\sin$, $\cos$ og $\tan$ — der kun afhænger af \emph{en af de to ikke-rette vinkler}. Det er det smukke: samme trekant, uanset størrelse, har samme forhold mellem siderne, hvis vinklen er den samme.}

\noindent Enhedscirklen er en cirkel med radius 1. For en vinkel $\theta$ målt fra den positive $x$-akse:
\begin{itemize}[leftmargin=2em,itemsep=2pt]
\item $\cos\theta$ er $x$-koordinaten for punktet på cirklen.
\item $\sin\theta$ er $y$-koordinaten for punktet på cirklen.
\item $\tan\theta=\dfrac{\sin\theta}{\cos\theta}$.
\end{itemize}

\regel{Definitioner (huskes på remse: \textbf{H}os \textbf{M}od $\Rightarrow$ \textbf{HM} for $\cos$ og $\sin$):}{%
\[
\sin\theta=\frac{\text{modstående}}{\text{hypotenuse}}, \qquad
\cos\theta=\frac{\text{hosliggende}}{\text{hypotenuse}}, \qquad
\tan\theta=\frac{\text{modstående}}{\text{hosliggende}}.
\]
}

\eksempel{En retvinklet trekant med vinkler 30, 60, 90 og hypotenuse 10}{%
  Modstående katete til 30$^\circ$: $10\cdot\sin 30^\circ=10\cdot 0{,}5=5$.\\[3pt]
  Hosliggende katete til 30$^\circ$: $10\cdot\cos 30^\circ=10\cdot 0{,}866\approx 8{,}66$.\\[3pt]
  \textit{Tjek:} $5^2+8{,}66^2=25+75=100=10^2$ — Pythagoras holder.
}

\eksempel{Stigen op ad en mur}{%
  En stige på 4\,m står op ad en lodret mur og danner vinklen 70$^\circ$ med jorden. Hvor højt op rør stigen muren?\\[3pt]
  Stigen er hypotenusen (4\,m). Den lodrette side er \textbf{modstående} til 70$^\circ$.\\[3pt]
  $\sin 70^\circ=\dfrac{\text{modstående}}{\text{hypotenuse}}=\dfrac{h}{4}$.\\[3pt]
  $h=4\cdot\sin 70^\circ\approx 4\cdot 0{,}9397\approx 3{,}76$\,m.\\[3pt]
  \textit{Tolkning:} stigen rør muren i ca.\ 3,76 meters højde.
}

\eksempel{Bestem vinklen i en trekant med kateter 3 og 5}{%
  Vi har en retvinklet trekant, den hosliggende katete (til vinklen) er 3, den modstående er 5.\\[3pt]
  $\tan\theta=\dfrac{5}{3}\approx 1{,}667$.\\[3pt]
  $\theta=\tan^{-1}(1{,}667)\approx 59{,}0^\circ$.\\[3pt]
  \textit{Tjek:} Pythagoras giver hypotenusen $\sqrt{3^2+5^2}=\sqrt{34}\approx 5{,}83$.
}

\laerertip{3.1--3.4: start med ren ``læs definitionen''-spørgsmål, slut med anvendelser. 3.4 (bonus) er den klassiske ``3-4-5''-familie (5-12-13 er Pythagoræisk). Eleven skal finde \emph{alle tre} vinkler, så de træner både $\sin^{-1}$ og $\cos^{-1}$.}

\opg{Opgave 3.1}{I en retvinklet trekant er den hosliggende katete 8 og hypotenusen 17. Beregn (a) $\cos\theta$, (b) $\sin\theta$, (c) $\tan\theta$.}
\facit{Hosliggende 8, hypotenuse 17.\\[2pt]
(a) $\cos\theta=\tfrac{8}{17}\approx 0{,}4706$.\\[2pt]
(b) Pythagoras: modstående $=\sqrt{17^2-8^2}=\sqrt{225}=15$. $\sin\theta=\tfrac{15}{17}\approx 0{,}8824$.\\[2pt]
(c) $\tan\theta=\tfrac{15}{8}=1{,}875$.}

\opg{Opgave 3.2}{En mand ser op på toppen af en bygning. Han står 30\,m fra bygningens fod, og hans synslinje danner vinklen 35$^\circ$ med vandret. Hvor høj er bygningen?}
\facit{Den lodrette side er modstående, den vandrette (30\,m) er hosliggende.\\[2pt]
$\tan 35^\circ=\tfrac{h}{30}\Rightarrow h=30\cdot\tan 35^\circ\approx 30\cdot 0{,}7002\approx 21{,}0$\,m.\\[2pt]
\textit{Tjek med $\sin$:} afstand til toppunkt $=\tfrac{30}{\cos 35^\circ}\approx 36{,}6$\,m, og $36{,}6\cdot\sin 35^\circ\approx 21{,}0$\,m — samme resultat.}

\opg{Opgave 3.3}{En flyvemaskine letter i en vinkel på 12$^\circ$ med vandret. Efter 500\,m langs flyveretningen, hvor højt er den? Tip: hypotenusen er 500\,m.}
\facit{Hypotenuse 500\,m, vinkel 12$^\circ$ med vandret. Den lodrette højde er modstående.\\[2pt]
$h=500\cdot\sin 12^\circ\approx 500\cdot 0{,}2079\approx 104{,}0$\,m.\\[2pt]
\textit{Bemærk:} det er en flad vinkel, så flyet er ikke så højt som man måske tror — kun ca.\ 100\,m op efter en halv kilometer.}

\opg{Opgave 3.4 (bonus)}{En retvinklet trekant har siderne 5, 12 og 13. Find alle tre spidse vinkler (i grader, afrundet til én decimal). Tip: start med den mindste vinkel.}
\facit{Sider: 5, 12, 13. Hypotenuse 13.\\[2pt]
Lille vinkel (modstående 5): $\sin\theta_1=\tfrac{5}{13}\Rightarrow\theta_1=\sin^{-1}(\tfrac{5}{13})\approx 22{,}6^\circ$.\\[2pt]
Stor vinkel (modstående 12): $\theta_2=\sin^{-1}(\tfrac{12}{13})\approx 67{,}4^\circ$.\\[2pt]
Tjek: $22{,}6^\circ+67{,}4^\circ+90^\circ=180^\circ$ — sum af vinkler i trekant.}

\end{document}
